\documentclass[../main]{subfiles}
\begin{document}
\chapter{Sistemas Termodinámicos:}
\img{images/02/tipos-sistemas-termodinamicos.png}{Sistemas termodinámicos.}
\info{
Contenido}{
Sistemas termodinámicos: sistemas abierto, cerrado, aislado, estados de equilibrio, diagrama de Clapeyron, trabajo interno, externo, trabajo de circulación.}


\actividad{ Leer, investigar y responder.}
\begin{enumerate}
    \item Definir sistemas termodinámicos, clasificarlos y explica cada uno.
    \item Definir componentes y fases.
    \item Definir tipos de interacciones.
    \item Definir estado y estado de equilibrio.
    \item Explicar las diferentes variables termodinámicas.
    \item Definir proceso termodinámico.
    \item Definir propiedad intensiva, extensiva.
    \item Explicar diferencia entre trabajo Externo e Interno, grafico y fórmulas signos.
    \item Explicar trabajo de circulación con diagrama de CLAPEYRON.
    
\end{enumerate}

\actividad{Resolver los siguientes ejercicios.}
\begin{enumerate}
    \item Determinar la cantidad de trabajo necesaria para que $1m^3$ de gas entre en un tanque que ya tiene el mismo gas a una presión de $5atm$.

\item Un cilindro con un pistón contiene un gas a una presión de $3atm$. Se desea introducir $0.5m^3$ de gas en el cilindro. ¿Cuánto trabajo se necesita realizar para lograrlo?

\item En una instalación de almacenamiento de gas, se tiene un tanque a una presión de $5atm$. Se requiere agregar $1m^3$ de gas al tanque. ¿Cuál es el trabajo necesario para llevar a cabo esta tarea?
\begin{center}
\includegraphics[scale=0.05]{images/02/tanque.jpg}
\end{center}

\item Un sistema de refrigeración utiliza un compresor para mantener la presión del gas a $2atm$ dentro de un contenedor. Si se desea introducir $2m^3$ de gas en el contenedor, ¿cuánto trabajo debe realizar el compresor?

\item En un laboratorio de investigación, se tiene un recipiente sellado que contiene gas a una presión de $4atm$. Para realizar un experimento, se necesita agregar $0.8m^3$ de gas al recipiente. ¿Cuál es el trabajo requerido para lograr esto?

\item En una planta de producción química, se cuenta con un reactor que opera a una presión de $6atm$. Se planea introducir $1.5m^3$ de gas en el reactor para iniciar una reacción. ¿Cuánto trabajo se debe realizar para completar esta tarea?

\img{images/02/esferagas.jpg}{Esferas de almacenamiento de gas.}

    \item Calcular el trabajo efectuado por el gas contenido en un cilindro se muestra con volumen inicial $v_i=0,025m^3$ y presión $p_i=2,5atm$, si al calentarle por su base $AB$ el volumen aumenta a $v_f=0,175m^3$ permaneciendo constante la presión.
    \begin{center}
        %\includegraphics[scale=0.05]{images/02/piston.jpg}
        \subfile{images/02/piston2}
    \end{center}
    \item Calcular el trabajo que es necesario efectuar sobre $v=0,5m^3$ de gas para que ingrese a un tanque que contiene el mismo gas a presión de $p=3atm$.
    \item Calcular el trabajo de circulación producido en el compresor de la figura si al mismo ingresa $v=1m^3$ de aire a la presión de la atmósfera y del mismo salen $v_f=0,025m^3$ a la presión de $p_f=3atm$, el trabajo de compresión ha sido $T=13500kgm$. 
    \begin{center}
        \includegraphics[scale=0.1]{images/02/compresor.jpg}
    \end{center}
 
\item Imagina que estás en una fábrica de producción de alimentos. En la línea de empaquetado, un compresor se utiliza para llenar bolsas de aire antes de sellarlas. El compresor toma $0.5m^3$ de aire a la presión atmosférica y lo comprime hasta un volumen de $0.02m^3$ a una presión de $3atm$. Durante este proceso de compresión, se realiza un trabajo de compresión de $8000J$.

¿Cuál es el trabajo de compresión producido en el compresor durante este proceso?


\end{enumerate}


\end{document}